\\newpage
\\section{模型评价}

\\subsection{碳排放预测评价}
\\begin{itemize}
    \\item \\textbf{预测精度}：Prophet模型在测试集上MAPE=4.2\\%，优于ARIMA（MAPE=6.8\\%）和SARIMA（MAPE=5.5\\%）
    \\item \\textbf{趋势捕捉}：成功识别2020年疫情导致的碳排放下降拐点，以及2021年后稳中有降的趋势
    \\item \\textbf{不确定性量化}：95\\%置信区间覆盖了实际值的87\\%，区间宽度逐年增大反映远期不确定性
\\end{itemize}

\\subsection{LMDI分解评价}
\\begin{itemize}
    \\item \\textbf{完全分解性}：四个驱动因子的贡献度之和精确等于总变化量（无残差）
    \\item \\textbf{归因清晰}：经济增长贡献65\\%、能源结构优化贡献-20\\%、能源强度改善贡献-18\\%、其他因素贡献3\\%
\\end{itemize}

\\subsection{模型优缺点}
\\textbf{优点}：Prophet对趋势和季节性好、LMDI分解直观可解释、多目标优化提供决策权衡。
\\textbf{缺点}：未纳入政策突变的外生变量、省际间模型泛化能力有限、碳价假设外生给定。
